\documentclass[12pt, a4paper]{article}

\usepackage[utf8]{inputenc}
\usepackage{geometry}
\geometry{margin=1in}

\title{Project Report: Web-Based Snake Game}
\author{Chetan & Vinay}
\date{29th of April, 2026}

\begin{document}

\maketitle

\begin{abstract}
This project is about building a Snake game that runs in a modern web browser using HTML5 Canvas , JavaScript and Python. The idea was to recreate the classic Snake game but also add some extra features like different types of food, increasing difficulty, and a temporary immunity feature. The game also connects to a backend to store scores.
\end{abstract}

\section{Introduction}
Snake game is one of the most popular classic games where the player controls a snake that keeps moving and grows longer as it eats food. The goal is to survive as long as possible without crashing.

In this project, I created a browser-based version of Snake. The main aim was to understand how games work using JavaScript, especially things like game loops, user input, and object-oriented programming.

\section{Gameplay and Mechanics}

\subsection{Controls and Setup}
When the game starts, the player is asked to enter their name and choose a difficulty level. The difficulty affects how fast the snake moves.

The controls are simple:
\begin{itemize}
    \item Arrow keys or WASD to move
    \item Spacebar to pause or resume the game
\end{itemize}

\subsection{Food System}
To make the game more interesting, I added different types of food:

\begin{itemize}
    \item \textbf{Normal Food:} Gives +1 score and increases the snake length.
    \item \textbf{Bonus Food:} Appears randomly and gives +5 score and increases the snake length.
    \item \textbf{Immune Food:} Appears after every few points. It gives temporary immunity where the snake can pass through walls.
\end{itemize}

There is also a small timer bar on the screen that shows how long the immunity will last.

\section{Technical Implementation}

\subsection{Graphics}
The game is drawn using the HTML5 canvas. JavaScript is used to update the snake position, draw the food, and refresh the screen continuously.

\subsection{Game Logic}
I used a simple object-oriented approach:
\begin{itemize}
    \item Snake class handles movement and body
    \item Food class handles spawning at random positions
\end{itemize}

The game runs using a loop that updates the state and redraws everything.

\subsection{Backend}
When the game ends, the score is sent to a backend server using a fetch request. The data includes the player's score and how the game ended.

\section{Conclusion}
Overall, this project helped me understand how games are built using JavaScript. Adding features like bonus food and wall immunity made the game more fun and less predictable.

\end{document}
